\section{Infinitesimal rigidity and finite observable coordinates}
\label{sec:v48-differential}

The exact single-offset inverse has a differential consequence which does
not follow from injectivity alone.  We prove it for variations of entire
analytic tables, allowing the contacts, obstacle shapes and marked lattice
to vary.  The analytic hypothesis is imposed on the parameter derivative
as well as on each table.  At the end of the section, a finite-dimensional
model is introduced only for the scalar-coordinate conclusion.

\subsection{Analytic variations and fixed intrinsic coordinates}

For $\rho>0$, let $\mathcal H_\rho$ be the real Banach space of bounded
$2\pi$-periodic holomorphic functions on $|\operatorname{Im}z|<\rho$
which are real on the real axis, with the supremum norm.  A
\emph{common-strip analytic-support family} is a map
\begin{equation}\label{eq:v48-support-family}
 t\longmapsto (h_{1,t},\ldots,h_{N,t},L_t)
       \in\mathcal H_\rho^N\times\mathrm{GL}^+(2,\mathbb R)
\end{equation}
which is $C^1$ on an open parameter interval and represents admissible
tables.  Families with a parameter in an open subset of $\mathbb R^p$
are defined in the same way.  Positive curvature, separation, and the
chosen clear-channel design persist near the base parameter.  A dot
means differentiation at that parameter.  Cauchy's estimates imply
$C^1$ dependence in each real $C^m$ support norm and holomorphy of
$\dot h_i$ on the strip.  Changes to moving contact frames preserve
these properties on every smaller strip.  No common bound for all
orders of differentiation is asserted.

Fix a table in $\mathscr C_N$ and a skeleton
$E=\mathcal T\cup\{f_1,f_2\}$ from
Theorem~\ref{thm:v45-clear-skeleton}.  Freeze all obstacle, channel and
deck marks, the transverse signs, and one sufficiently small positive
offset $d_e$ for each edge.  The same offset may be used on all edges.
In the smoothly moving intrinsic frame of each channel, its contacts
have coordinates $(0,0)$ and $(g_e,0)$ and the transverse variable is
$u$.  Fix contact-centered squares $Q_{e,b}=(-r_{e,b},r_{e,b})^2$ whose
closures lie in the respective positive-density regions, and nonzero scalar
anchors $a_{e,b}\in(-r_{e,b},r_{e,b})$.  These choices persist on a
parameter neighborhood.  Write
\[
       f_{e,b,t}=\frac{dP_{e,b,t}}{du\,dv}
\]
for the limiting conditional density in this frame.  The frame is a
coordinate convention for the intrinsic experiment, not an observed
inter-channel pose.

A common Euclidean gauge can be fixed by putting the support center of
a root obstacle at the origin and the oriented axis of a root channel
along the first coordinate axis.  The support center is the vector
$s(K)=\pi^{-1}\int_0^{2\pi}h_K(\theta)n(\theta)\,d\theta$.
The root channel has positive length; its direction is smooth.  This
gives a local $C^1$ slice for the three-dimensional proper Euclidean
action.  When needed, restrict the strip after this change of gauge.
All conclusions about zero tangent vectors in this slice mean equality
of the support-function and lattice derivatives, not a particular
parametrization of the boundary curves.

\begin{lemma}[Differentiability on an interior square]
\label{lem:v48-forward-C1}
For a family \eqref{eq:v48-support-family}, each gap, intrinsic contact
frame and finite contact jet is $C^1$.  On every fixed positive-density
square, $t\mapsto f_{e,b,t}$ is $C^1$ in each prescribed finite $C^m$
norm, after restricting the parameter neighborhood as necessary.
Consequently
\begin{equation}\label{eq:v48-test-derivative}
 \left.\frac{d}{dt}\right|_{0}\int\phi\,dP_{e,b,t}
       =\int_{Q_{e,b}}\phi(u,v)\dot f_{e,b}(u,v)\,du\,dv,
       \qquad \phi\in C_c^\infty(Q_{e,b}).
\end{equation}
The derivatives in this statement are at fixed excess time $d_e$ in
the moving intrinsic transverse coordinates.
\end{lemma}

\begin{proof}
For a facing pair, the two endpoint stationarity equations have an
invertible Hessian.  In contact transverse coordinates the quadratic
part of the distance has Hessian
\[
 \begin{pmatrix}\kappa_0+g^{-1}&-g^{-1}\\
                 -g^{-1}&\kappa_1+g^{-1}\end{pmatrix},
 \qquad
 \det=\kappa_0\kappa_1+g^{-1}(\kappa_0+\kappa_1)>0.
\]
The support parametrization and its finite derivatives are $C^1$ by
Cauchy's estimates.  The implicit function theorem therefore gives
$C^1$ contact points, gap and frame; it also gives $C^1$ local graph
functions at every finite smooth order.  Clearance fixes the selected
facing pair and its gates on this neighborhood.

Apply the weighted half-line construction and the differentiated
relative law to this family.  The contraction constants are uniform
on a sufficiently small neighborhood with positive gap and curvature
margins.  Differentiating the fixed-point equation solves for the
parameter derivative with the same bounded inverse; finite endpoint
derivatives use the higher smooth bounds already supplied by
\eqref{eq:v48-support-family}.  The summable endpoint and terminal
estimates of Lemmas~\ref{lem:v22-weighted-inverse} and
\ref{lem:v22-envelope} justify the corresponding action derivatives.
For the amplitude, we make the functional-family derivative explicit.
Suppress the endpoint type and write
\[
 H_t^0=H_t(0),\qquad G_t=(H_t^0)^{-1},\qquad
 \Delta H_t(u)=H_t(u)-H_t^0,\qquad
 \mathcal T_t(u)=G_t\Delta H_t(u).
\]
The uniformly positive tridiagonal reference operator is $C^1$ in
operator norm.  The resolvent identity gives
\begin{equation}\label{eq:v49-moving-reference}
 \dot G=-G\dot H^0G,\qquad
 \dot{\Delta H}=\dot H(u)-\dot H^0,\qquad
 \dot{\mathcal T}=\dot G\,\Delta H+G\dot{\Delta H}.
\end{equation}
In particular the derivative of the moving reference is not omitted.
For each fixed endpoint order $m$, locality and the weighted half-line
estimates imply, on a uniform small neighborhood,
\[
 \sum_{i,j\ge1}\bigl|
       \partial_u^r\partial_t^a(\Delta H_t(u))_{ij}\bigr|
       \le C_m\sum_{i\ge1}\rho^i<\infty,
          \qquad 0\le r\le m,\quad a=0,1.
\]
Only neighboring entries occur.  The terms without endpoint decay in
$\dot H(u)$ cancel those in $\dot H^0$; the remaining terms contain
an orbit coordinate or its parameter derivative.  Fixed index powers
are absorbed in the strict exponential margin.  Finite entry sums
therefore converge with their derivatives uniformly in trace norm.
Thus $\mathcal T_t$ is $C^1$ in trace norm, in each fixed endpoint norm,
and $\|\mathcal T_t\|\le q<1$ after the collar is shrunk once.

The edge sum $U_t$ of \eqref{eq:v4-half-product} has the same
summability: $g_t[-\ell_{t,uv}(0,0)]=1$ identically in $t$, so its
reference derivative cancels before the edges are summed.  In the
absolutely convergent trace series \eqref{eq:v4-half-amplitude},
cyclicity and the trace-ideal inequality give
\[
 \frac{d}{dt}\operatorname{tr}(\mathcal T_t^n)
    =n\operatorname{tr}(\mathcal T_t^{n-1}\dot{\mathcal T}_t),
 \qquad
 |\operatorname{tr}(\mathcal T_t^{n-1}\dot{\mathcal T}_t)|
     \le q^{n-1}\|\dot{\mathcal T}_t\|_1.
\]
Summing this geometric majorant proves
\begin{equation}\label{eq:v49-amplitude-derivative}
 \frac{\dot B}{B}=\dot U-
       \operatorname{tr}\bigl((I+\mathcal T)^{-1}
                                      \dot{\mathcal T}\bigr).
\end{equation}
Fixed endpoint differentiation introduces only fixed polynomial factors
in the series index and bounded resolvent factors, which remain
summable.  This proves the amplitude's $C^1$ dependence at each finite
smooth order, directly from the same relative construction.  The
half-line actions have the $C^1$ dependence already established above.

It remains to check the normalization, since its integral meets a
moving boundary.  In a fixed box $U$ containing the selected supports,
put $A_t(u,v)=S_t(u)+S_t(v)$ and $w_t(u,v)=B_t(u)B_t(v)$.
The normalizer is
\[
                  Z_t=\int_U w_t(d_e-A_t)_+.
\]
The level set $A_t=d_e$ has two-dimensional measure zero.  Indeed
$S_t'$ vanishes only at zero, whereas $d_e>0$; the gradient of $A_t$
is nonzero at every point of this level set.  The positive-part map is
Lipschitz.  Its difference quotients are bounded by a common bound for
the parameter derivative of $A_t$.  Dominated convergence gives
\[
 \dot Z=\int_U\{\dot w(d_e-A)_+-w\dot A\,
                                     \mathbf1_{\{A<d_e\}}\}.
\]
The same argument as the parameter varies gives continuity of this
derivative: the indicators converge almost everywhere off the limiting
level set.  The positive lower bound for $Z_t$ persists.  On a positive
interior square the cutoff is inactive, so the product and quotient
rules prove the asserted $C^1(C^m)$ statement.  Integrating against the
fixed compactly supported test function proves
\eqref{eq:v48-test-derivative}.  In particular no derivative of a Dirac
mass on the cap boundary is introduced into this argument.
\end{proof}

\subsection{Linearization of the signed contact inverse}

\begin{lemma}[A zero law derivative has zero contact-jet derivative]
\label{lem:v48-local-kernel}
In a common-strip analytic-support family of one channel, suppose
$\dot g=0$ and $\dot f_b=0$ on the chosen interior square for both
same-type laws at their fixed known positive offset.  Then every
labelled anchored graph coefficient has zero derivative.  This holds
for odd as well as even orders and does not assume a known amplitude.
\end{lemma}

\begin{proof}
We display the differential of the amplitude-free inverse.  For one
type omit $b$, put $q=\sqrt{1-R_f(a,a)}>0$, and write
$t(u)=(1-R_f(u,a))/q$.  At a fixed positive density,
\begin{align}
 \dot R_f(u,v)&=R_f(u,v)\left\{
    \frac{\dot f(u,v)}{f(u,v)}+\frac{\dot f(0,0)}{f(0,0)}
   -\frac{\dot f(u,0)}{f(u,0)}-\frac{\dot f(0,v)}{f(0,v)}\right\},
                         \label{eq:v48-ratio-differential}\\
 \dot t(u)&=-\frac{\dot R_f(u,a)}q
                   +\frac{t(u)\dot R_f(a,a)}{2q^2},\qquad
 \dot S(u)=\frac{d_e\dot t(u)}{(1+t(u))^2}.
                         \label{eq:v48-action-differential}
\end{align}
All these operations are $C^1$ in any prescribed finite $C^m$ norm:
evaluation on a fixed slice is bounded linear, the denominators are
positive, and the sole square root is a scalar evaluated at the nonzero
anchor.  Lemma~\ref{lem:v48-forward-C1} and
Theorem~\ref{thm:v26-density-inverse} allow the ordinary chain rule.
Thus $\dot f=0$ implies $\dot S=0$, including every finite derivative
at zero.  The two signed coordinates have not been identified.

Write $a_b=S_b''(0)$.  The explicit leading recovery is
\[
 \gamma=\operatorname{arsinh}(g\sqrt{a_0a_1}),\qquad
 c_b=\cosh\gamma\sqrt{a_b/a_{1-b}},\qquad
 \kappa_b=(c_b-1)/g.
\]
It follows that $\dot\kappa_b=0$.  At order $n\ge3$, use the actual
finite-jet identity
\[
             s_n=M_nq_n+R_n,
\]
where $R_n$ depends only on lower graph jets and the gap, and $M_n$ is
the invertible block in \eqref{eq:v22-general-jet-block}.  Its validity
for smooth representatives, including dependence on finite rather than
formal jets, is supplied by
Lemma~\ref{lem:v27-smooth-jet-factorization}.  Each finite map is smooth,
as also follows from the finite Bellman recursion.  Differentiating gives
\begin{equation}\label{eq:v48-jet-differential}
       \dot q_n=M_n^{-1}\{\dot s_n-\dot R_n-\dot M_nq_n\}.
\end{equation}
Induction, beginning with the leading recovery, makes every term on
the right zero.  The orders zero and one are fixed by the contact
anchors.  The induction is carried out separately at each finite
$n$; it neither differentiates an infinite inverse series nor uses an
order-independent bound on $M_n^{-1}$.
\end{proof}

\begin{lemma}[Analytic propagation of a stationary contact variation]
\label{lem:v48-analytic-variation}
If all anchored contact-graph coefficients of a common-strip
analytic-support family have zero derivative and $\dot g=0$, the entire
incident obstacle images are stationary to first order in that channel
frame.  Their support functions therefore have zero derivative in every
real $C^m$ norm.
\end{lemma}

\begin{proof}
Express the two bodies in the same moving channel frame.  The contact
origins are $0$ and $(g,0)$, whose derivatives vanish under the hypothesis.
The normal directions at the contacts are fixed, and any reversal
needed to use the positively oriented tangent is fixed as well.
The local graph-to-support conversion is a smooth triangular map at
each finite order, with positive curvature as its nondegeneracy
condition; see \eqref{eq:v46-graph-support}.  Hence every derivative in
normal angle of the support variation vanishes at the contact normal.

For clarity, the analyticity needed at this step concerns the variation
itself.  If $A_t x=R_{\alpha_t}x+a_t$ is the placement of a channel frame,
the support of a body in that frame is obtained from its global support
by an angular translation and a linear first-harmonic term.  Its
parameter derivative is a sum of the translated $\dot h$, a multiple
of the translated $h'$, and first harmonics.  Each is holomorphic on
any fixed smaller strip.  This follows from
\eqref{eq:v48-support-family}, Cauchy's estimates and the $C^1$
frame dependence, not from pointwise analyticity of $h_t$ alone.
The identity theorem now makes the support variation zero on the
connected strip, hence on the entire real normal circle.  Its real
$C^m$ norm is zero for each $m$.  No norm-continuous inverse for
analytic continuation is used or concluded.
\end{proof}

\subsection{Differentiable incidence matching and the lattice}

\begin{lemma}[Finite harmonic registration on the asymmetric locus]
\label{lem:v48-registration}
For each obstacle with trivial proper Euclidean symmetry, choose a
finite set $F\subset\{2,3,\ldots\}$ of nonzero support harmonics whose
indices have greatest common divisor one.  These choices persist
locally.  Incidence congruences, the rotations of all selected channel
frames into one common frame, the lattice $L$, and representative obstacle
placements are $C^1$ functions of the complete channel-frame images
along common-strip analytic-support families.  If all these images are
stationary to first order, the reconstructed gauge-fixed table is
stationary to first order.
For noncircular obstacles, the same conclusions hold on a local alignment
branch through the actual base realization, without asserting globally
single-valued registration.
\end{lemma}

\begin{proof}
The finite witness exists by
Proposition~\ref{prop:v45-generic-asymmetry}: otherwise the nonzero
centered support harmonics have a common divisor larger than one, or
there are none, giving respectively a nontrivial rotation or a circle.
The greatest common divisors of successive finite subsets stabilize,
so a finite witness suffices when the full divisor is one.  Nonzero
coefficients remain nonzero on a neighborhood.

Choose integers $a_k$ with $\sum_{k\in F}a_kk=1$.  For two copies
$C_s,C_t$ of this obstacle, with $C_t=UC_s+a$ and
$U=e^{i\alpha}$, define
\begin{equation}\label{eq:v48-bezout-registration}
 r_k=\frac{z_k(C_t)}{z_k(C_s)},\qquad
 U=\prod_{k\in F}r_k^{-a_k},\qquad
 a=s(C_t)-Us(C_s),
 \quad z_k(C)=\frac1{2\pi}\int_0^{2\pi}h_C(\theta)e^{-ik\theta}\,d\theta.
\end{equation}
Indeed $r_k=e^{-ik\alpha}$, so the product is $e^{i\alpha}$.
Only integer powers of nonzero scalars occur; no argument or root
branch is chosen.  Centers and coefficients are bounded linear
functionals of the real support function.  Formula
\eqref{eq:v48-bezout-registration} is consequently $C^1$ on its
nonvanishing domain.  It recovers the unique proper incidence
congruence.  For $F=\{2,3\}$ it reduces to
\eqref{eq:v45-general-registration}.  Thus the present argument does
not restrict the asymmetric locus to the second-and-third-harmonic
subclass.

Propagate these congruences along the finite spanning tree from a fixed
root incidence.  Match the initial obstacle of either remaining edge
to get its rotation as well.  This is a finite composition of $C^1$
operations.  In the resulting orientation let $R_e$ be the rotation of
edge $e$ and set
\[
 d_e^{\rm cen}=R_e\{s(C_{e,+})-s(C_{e,-})\},\qquad
 p_1=0,\qquad p_a=\sum_{e\in[1,a]_{\mathcal T}}\epsilon_e d_e^{\rm cen}.
\]
Here $m_a$ denotes the sum of signed deck marks along the same tree
path.  For $f_i=(a_i,b_i,\ell_i)$ the marked gain is
$\eta_i=m_{a_i}+\ell_i-m_{b_i}$ and
\begin{equation}\label{eq:v48-differential-cochain}
 v_i=p_{a_i}+d_{f_i}^{\rm cen}-p_{b_i},\qquad
 L=(v_1\ v_2)M^{-1},\qquad c_a=p_a-Lm_a,
 \qquad M=(\eta_1\ \eta_2).
\end{equation}
These are the displacement-cochain identities of
Proposition~\ref{prop:v45-explicit-image-inverse}, whose telescoping
proof depends on the recovered rotations, not on using harmonics two
and three in particular.  The matrix $M$ is fixed and invertible; it
need not be unimodular.  Centering and rotating one recovered copy of
each obstacle, then translating it to $c_a$, completes the
reconstruction in the root-center gauge.  All operations are $C^1$
along the family.  In particular
\[
                  \dot L=(\dot v_1\ \dot v_2)M^{-1}.
\]
If every input support variation is zero, every coefficient and center
derivative, every relative-rotation derivative, and then every quantity
in \eqref{eq:v48-differential-cochain} has zero derivative.  The output
support variations and lattice variation are zero.  This proves the
last assertion without assuming a differentiable inverse to analytic
continuation.

\paragraph{Local registration at a noncircular table.}
When the base obstacle has a nontrivial finite symmetry group, the
Bezout formula in Lemma~\ref{lem:v48-registration} need not exist.
Use instead the local angular lift of
Lemma~\ref{lem:v49-local-alignment}, separately at each link of a fixed
channel-incidence spanning tree.  Such branches exist because the actual
continuously moving channel frames give congruences between each fixed
pair of labelled occurrences.  Every noncircular base obstacle has
one nonzero harmonic, so these local choices persist along a
$C^1$ family, whether or not its finite symmetry is preserved.
Stationary channel-frame support images give stationary relative
rotations and translations on these branches.  All subsequent center
and cochain operations in \eqref{eq:v48-differential-cochain} are the
same finite $C^1$ operations.  Hence stationary images again give
$\dot L=0$ and stationary representative supports in the common gauge.
The actual base placements are used only to specify the local branch
for this derivative argument, not supplied as extra observations.
The finite alternatives in Theorem~\ref{thm:v49-finite-fiber} are not
removed globally by this local choice.
\end{proof}

\begin{theorem}[Infinitesimal single-offset periodic rigidity]
\label{thm:v48-infinitesimal}
Let $T_t$ be a common-strip analytic-support family through
$T_0\in\mathscr C_N$, with a persistent selected skeleton and fixed
positive offsets as above.  The following are equivalent:
\begin{enumerate}
\item $\dot g_e=0$ and $\dot f_{e,b}=0$ on $Q_{e,b}$ for every $e,b$;
\item $\dot g_e=0$ and $\left.\frac{d}{dt}\right|_0
\int\phi\,dP_{e,b,t}=0$ for every $e,b$ and every real
$\phi\in C_c^\infty(Q_{e,b})$;
\item the table derivative is generated by one common infinitesimal
proper Euclidean motion.
\end{enumerate}
In particular the data derivative is injective in the root-center and
root-direction gauge, including variations of the marked lattice.
\end{theorem}

\begin{proof}
The first two assertions are equivalent by
Lemma~\ref{lem:v48-forward-C1}.  For the nontrivial direction, a
continuous nonzero $\dot f$ has constant strict sign on some small
open ball.  A nonnegative smooth function supported in that ball has
nonzero integral against $\dot f$, contradicting the second assertion.

Assume the first assertion.  Lemma~\ref{lem:v48-local-kernel} makes all
contact-jet variations zero.  Lemma~\ref{lem:v48-analytic-variation}
makes all entire channel-frame support variations zero.
Lemma~\ref{lem:v48-registration} then makes the gauge-fixed table
variation zero.  The $C^1$ gauge change therefore expresses the
original derivative as a common infinitesimal translation and rotation.
Explicitly, for constants $a\in\mathbb R^2$ and $\omega\in\mathbb R$,
it has the support and lattice form
\begin{equation}\label{eq:v48-gauge-kernel}
       \dot h_i(\theta)=a\cdot n(\theta)-\omega h_i'(\theta),
                   \qquad \dot L=\omega J L,
       \qquad J=\begin{pmatrix}0&-1\\1&0\end{pmatrix}.
\end{equation}
The converse follows from invariance of gaps and signed intrinsic
conditional laws under a simultaneous proper Euclidean motion.  All
parameter derivatives are ordinary derivatives along the stipulated
family.  Exact injectivity of the undifferentiated map is not used as
a substitute for any step of the kernel argument.
\end{proof}

\subsection{Finite scalar coordinates on immersed models}

\begin{theorem}[Local coordinates from finitely many gap and law tests]
\label{thm:v48-finite-coordinates}
Let $\Theta\subset\mathbb R^p$ be open, $p\ge1$, and let
$\theta\mapsto T_\theta$ be a $C^1$ common-strip analytic-support family
in the fixed Euclidean gauge.  Suppose its table derivative at
$\theta_0$ is injective and $T_{\theta_0}\in\mathscr C_N$.
Use a persistent skeleton and fixed positive offsets.  Then there are
$p$ scalar observables $\lambda_1,\ldots,\lambda_p$, each either
$g_e$ or $\int\phi\,dP_{e,b}$ for a real
$\phi\in C_c^\infty(Q_{e,b})$, such that
\begin{equation}\label{eq:v48-finite-coordinate-map}
 F(\theta)=(\lambda_1(T_\theta),\ldots,\lambda_p(T_\theta))
\end{equation}
is a $C^1$ diffeomorphism between neighborhoods of $\theta_0$ and
$F(\theta_0)$.  On a sufficiently small convex parameter ball there
are constants $c,C>0$ such that
\begin{equation}\label{eq:v48-local-bilip}
 c\,|\theta-\widetilde\theta|
       \le |F(\theta)-F(\widetilde\theta)|
       \le C\,|\theta-\widetilde\theta|.
\end{equation}
For each fixed $m$, after restricting this ball if necessary, the
$C^m$ differences of the gauge-fixed support functions and the
matrix difference of $L$ are bounded by a constant times
$|F(\theta)-F(\widetilde\theta)|$.
\end{theorem}

\begin{proof}
Consider on $\mathbb R^p$ the linear functionals
\[
 v\longmapsto Dg_e(\theta_0)v,\qquad
 v\longmapsto\int_{Q_{e,b}}\phi\,Df_{e,b}(\theta_0)v,
                 \qquad \phi\in C_c^\infty(Q_{e,b}).
\]
Their common kernel is zero.  Indeed a vector in that kernel gives a
table variation with vanishing intrinsic derivative by
Theorem~\ref{thm:v48-infinitesimal}; the gauge and the immersion
hypothesis then make the vector zero.  A set of linear functionals
with zero common kernel spans the dual of a finite-dimensional vector
space: a proper span would have a nonzero annihilator.  Select a
basis of $p$ members from this set.  Each selected member is the
derivative of an allowed scalar observable, not a derivative evaluation
of an unknown density.  Lemma~\ref{lem:v48-forward-C1} makes the
resulting map $F$ continuously differentiable.  Its derivative
$A=DF(\theta_0)$ is invertible, so the inverse function theorem proves
the local coordinate assertion.

There is also a direct bound avoiding any assertion that pointwise
nonsingularity alone implies a lower Lipschitz estimate.  On a small
convex ball, continuity gives
$\|DF(\theta)-A\|\le(2\|A^{-1}\|)^{-1}$.  For
$z=\theta-\widetilde\theta$, integration along the joining segment
therefore yields
\[
 |F(\theta)-F(\widetilde\theta)|
 \ge |Az|-\frac{|z|}{2\|A^{-1}\|}
 \ge\frac{|z|}{2\|A^{-1}\|}.
\]
The upper bound follows from the local bound for $DF$.  Finally the
$C^1(\mathcal H_\rho)$ family has locally bounded first parameter
derivatives in every fixed real $C^m$ support norm, by Cauchy's
estimates on a smaller strip, and bounded lattice derivatives.
Integrating these derivatives on the same ball and using the lower
bound proves the last assertion.
\end{proof}

Theorem~\ref{thm:v48-finite-coordinates} is a model-local coordinate
result for exact expectations.  It does not claim a fixed finite
observation vector for all analytic tables, a dimension bound for an
arbitrary compact analytic class, or a statistical efficiency theorem.
The full-class finite-histogram reconstruction below uses increasing
resolution and its separately stated conditioning hypotheses.  The
fixed-contact variation bundle and uniform positive designs in
Section~\ref{sec:v25-analytic-variation-bundle} concern a different question:
realizing compatible finite jets and detecting boundary velocities with
a finite collection of offsets.  Here the gaps, contacts and lattice
may move, and a single offset per channel suffices for the differential
law inverse.  The analytic identity theorem removes the shape kernel;
finite harmonic registration removes the relative-frame and lattice
kernel; finite-dimensionality is used only in the final selection of
scalar tests.
